\naslov{Linearna NDE prvega reda}

\pojem{Linearna NDE prvega reda} je enačba oblike
\[
  y' = f(x) y + g(x),
\]
kjer sta $f(x)$ in $g(x)$ znani funkciji.
Ta enačba je \pojem{nehomogena} z \pojem{nehomogenostjo} $g(x)$.
Njena \pojem{homogenizacija} je enačba
\[
  y' = f(x) y.
\]
Predpostavimo, da je $y(x) \in \zvodv{1}{[a, b]}$.
Oglejmo si operator $A : \zvodv{1}{[a, b]} \to \zvezne{[a, b]}$,
definiran kot
\[
  A y(x) = y'(x) - f(x) y.
\]

\begin{trditev}
  Preslikava $A$ je linearen operator.
\end{trditev}
Dokaz je trivialen, in zato izpuščen.
Vidimo, da je $y$ rešitev homogene enačbe natanko tedaj, ko je $A(y) = 0$.
Rešitev homogene enačbe je torej jedro preslikave $A$.
Homogena enačba je enačba z ločljivimi spremenljivkami, torej jo znamo rešiti.
Rešitve so oblike
\[
  y(x) = C \exp\left( \int_a^x f(\xi) d\xi \right)
\]
za $C \in \R$.
Množico teh rešitev označimo z $R_h$.

\begin{trditev}
  Naj bosta $y_1, y_2$ rešitvi nehomogene enačbe.
  Tedaj je $y(x) = y_1(x) - y_2(x)$ rešitev homogene enačbe.
\end{trditev}

\begin{proof}
  Izračun odvoda nam da
  \[
	(y_1(x) - y_2(x))' = f(x) y_1(x) + g(x) - (f(x) y_2(x) + g(x)) = f(x)
	(y_1(x) - y_2(x)).
  \]
\end{proof}

\begin{definicija}
  Naj bo $V$ nek vektorski prostor in $W \subseteq V$.
  Če obstaja tak vektorski podprostor $H \subseteq V$, da za poljubna $w_1, w_2
  \in W$ velja $w_1 - w_2 \in H$, je $W$ \pojem{afin podprostor} v $V$,
  modeliran z vektorskim podprostorom $H$.
\end{definicija}

Rešitve nehomogene enačbe so torej afin prostor, modeliran s prostorom $R_h$
rešitev homogene enačbe.
Če želimo poiskati splošno rešitev, poiščemo rešitev homogenega sistema, in neko
partikularno rešitev.
Partikularno rešitev dobimo z nastavkom
\[
  y_p(x) = C(x) \exp\left( \int_a^x f(\xi) d\xi \right),
\]
temu postopku pravimo \pojem{variacija konstante}.

\vprasanje{Kako rešiš linearno NDE prvega reda? Utemelji postopek.}

S tem znanjem lahko rešimo še dve posebni NDE.
Prva je Bernoulijeva enačba
\[
  p(x) y' + q(x) y = r(x) y^\alpha(x)
\]
za $\alpha \in \R$.
V primeru $\alpha = 0$ ali $\alpha = 1$, je to nehomogena linearna enačba prvega
reda.
Sicer vpeljemo $z(x) = (y(x))^{1-\alpha}$ in računamo
\[
  p(x) z'(x) \inv{1-\alpha} + q(x) z(x) = r(x)
\]
oziroma
\[
  z'(x) + \frac{q(x)}{p(x)} (1-\alpha) z(x) = \frac{r(x)}{p(x)} (1-\alpha).
\]
To je nehomogena linearna NDE prvega reda, torej jo znamo rešiti.

\vprasanje{Kako rešiš Bernoulijevo enačbo?}

Druga taka enačba je Riccatijeva enačba
\[
  y'(x) = a(x) y^2(x) + b(x) y(x) + c(x),
\]
ki je v splošnem ne znamo rešiti.
Poznamo pa dva načina obravnave, ki nas lahko včasih pripeljeta do rešitve.
Denimo, da uganemo neko partikularno rešitev $y_p(x)$.
Enačbo tedaj rešujemo z nastavkom $y(x) = y_p(x) + z(x)$ za neko neznano
funkcijo $z$.
Če to vstavimo v enačbo, dobimo
\[
  y_p' + z' = a y_p^2 + 2 a y_p z + a z^2 + b y_p + bz + c,
\]
členi $y_p'$, $a y_p^2$, $b y_p$ in $c$ odpadejo, ker tvorijo rešitev enačbe.
Ostane torej
\[
  z' = (2a y_p + b) z + a z^2,
\]
kar je Bernoulijeva enačba, ki jo znamo rešiti.

Drug način za reševanje Riccatijeve enačbe je s pretvorbo na linearni sistem
prvega reda.
Vpeljemo $y = \nicefrac{u}{v}$, s čimer dobimo
\[
  u'v - uv' = au^2 + buv + cv^2.
\]
Ker imamo dve neznanki, potrebujemo še eno enačbo.
Izberemo $u'v = buv + cv^2$.
Iz tega izpeljemo $v' = - a u$ in $u' = bu + cv$.
Zapisano matrično
\[
  \begin{bmatrix}
	v' \\ u'
  \end{bmatrix}
  =
  \begin{bmatrix}
	0 & -a \\
	c & b
  \end{bmatrix}
  \begin{bmatrix}
	v \\ u
  \end{bmatrix}
\]
Sistema v splošnem ne znamo rešiti, ker funkcije $a, b, c$ niso konstantne.
Lahko pa rešitev zapišemo v obliki neskončne vrste.

\vprasanje{Kako rešiš Riccatijevo enačbo?}